\documentclass[12pt]{article}
\usepackage[utf8]{inputenc}
\usepackage[T1]{fontenc}
\usepackage{hyperref}
\title{Sowa Language Resources for K-12 Curriculum Development}
\author{Compiled by Claude Sonnet 5.5 (claude-sonnet-5-5) for Andr\'as Kornai}
\date{}
\begin{document}
\maketitle

\section{Introduction}
Sowa (ISO 639-3 sww) is a dormant Austronesian language of south-central Pentecost Island, Vanuatu. Search results report that it was displaced by the neighbouring language Apma in the 20th century and that its last native speaker, Maurice Tabi, died in 2000. Very little documentation exists. No curriculum-ready materials were found.

\section{Primary Educational Resources}
No primers, textbooks or school materials in Sowa were found in search results.

\section{Linguistic and Reference Materials}
\begin{itemize}
\item Wikipedia overview of the Sowa language: \url{https://en.wikipedia.org/wiki/Sowa_language}
\item The search summary reports that David Walsh collected a vocabulary list in 1969. I did not locate the list itself.
\item Cynthia Schneider, \emph{A grammar of Abma} (2010), describes the neighbouring language that replaced Sowa, not Sowa itself: \url{https://tlopo.clld.org/sources/Schneider2010}
\end{itemize}

\section{Religious and Community Resources}
None found.

\section{Digital and Audio Resources}
\begin{itemize}
\item PARADISEC catalogue entry for Sowa, which per the search summary records only an archived web page by Andrew Gray: \url{https://catalog.paradisec.org.au/repository/External/Sowa}
\end{itemize}

\section{Curriculum Development Context}
Sowa has no first-language speakers and is not known to be taught in schools. A revitalization curriculum would have to be reconstructive. It would start from the 1969 word list and the archived material, and it would draw on related languages of Pentecost such as Apma, with community involvement. Little exists, and it should not be overstated.
\end{document}